\Def Решето Эратосфена - алгоритм, позволяющий найти все простые числа от 2 до n.

\section*{Решето Эратосфена}

\Alg 
\begin{enumerate}
    \item Берем минимальное непросмотренное число k, объявляем его простым.
    \item Далее все числа от 2 до n, кратные k объявляются непростыми.
    \item Повторить пока все числа не рассмотрены.
\end{enumerate}

\begin{lstlisting}
vector <bool> primes(n, true);  
primes[0] = primes[1] = false;
for (int i = 2; i < n; ++i) {
  if (!primes[i]) {
    continue;
  }
  for (int j = 2; i * j < n; ++j) {  
    primes[i * j] = false;  
  }
}
\end{lstlisting}

\textbf{Оптимизация.}
\begin{enumerate}
    \item Можно увеличивать $i$ на 2 каждый раз, так как кроме двойки все простые нечетные.
    \item Можно начинать перебирать $j$ начиная с $i$ , поскольку все числа меньшие $i^2$, кратные $i,$ обязательно имеют простой делитель меньше $i$, а значит, все они уже были отсеяны раньше.
\end{enumerate}

\textbf{Сложность решета Эратосфена.}


\Thbd Сложность алгоритма выше составляет $\mathcal{O}(n\log\log n)$, где $n$ - число, до которого ищем простые.

\section*{Линейное решето Эратосфена}

\Alg 
\begin{enumerate}
    \item Для каждого числа от 2 до n найдем минимальный простой делитель (массив mind). Тогда простые легко детектируются.
    \item При рассмотрении очередного числа i пройдемся по массиву уже найденных простых primes и будем обновлять минимальный делитель у просматриваемых чисел.
\end{enumerate}

\pagebreak

\begin{lstlisting}
vector<int> primes;
vector<int> mind(n + 1);

for (int i = 2; i <= n; ++i) {
  mind[i] = i;
}

for (int i = 2; i <= n; ++i) {
  if (mind[i] == i) {
    primes.push_back(i);
  }

  for (int p : primes) {
    if (p * i > n || p > mind[i]) {
      break;
    }
    mind[i * p] = p;
  }
}
\end{lstlisting}

\textbf{Сложность линейного решета Эратосфена.}

\begin{itemize}
    \item Внешний цикл for делает $O(n)$ операций: очевидно.
    \item Внутренний цикл также делает суммарно $O(n)$ операций: докажем.

    \Proof\\
    \textbf{Лемма 1}: $\forall n \in \mathbb{N}\;\; \exists!$ разложение $n = a \cdot b$, где $a$ - наименьший простой делитель $n$, а $b$ - остальная часть. Единственность очевидна из основной теоремы арифметики.\\
    \textbf{Лемма 2}: Каждая пара $(a, b)$, где $a$ - простое число, $b$ - натуральное число с наименьшим простым делителем $ > a$, делает ровно 1 изменение в массиве mind. Это следует из Леммы 1 (Каждая пара "бьёт"\ ровно в 1-у ячейку mind).\\
    \textbf{Лемма 3}: У нас не "сгенерируется"\ $> 1$ одинаковой пары $(a, b)$, то есть не будет такой ситуации, что $\exists a, b, c, d: (a, b), (c, d) \text{ и при этом } a \cdot b = c \cdot d$\\
    По итогу, так как у нас "генерируются"\ различные пары (a, b), каждая из которых обновляет ровно одну ячейку mind, то таких пар O(n) (Ведь mind.size() = O(n)), а значит вложенный массив суммарно делает O(n) операций.\\
    \Endproof
\end{itemize}

% Сложность определяется суммарным числом итераций внутреннего for.

% На первый взгляд для каждого p в разложении k изменяется mind[k]. Однако каждое mind[k] рассматривается ровно 1 раз, следовательно сложность линейная.

\pagebreak